\documentclass[11pt,a4paper]{article}
\usepackage[margin=1.7cm]{geometry}
\usepackage{fontspec}\usepackage{xeCJK}
\setCJKmainfont{Songti TC}\setCJKsansfont{PingFang TC}
\usepackage{graphicx,xcolor,fancyhdr}
\definecolor{red}{RGB}{210,30,30}
\pagestyle{fancy}\fancyhf{}\renewcommand{\headrulewidth}{0pt}\fancyfoot[C]{\small\thepage}
\setlength{\parindent}{0pt}\setlength{\parskip}{0pt}
% \step{number}{image}{max height}{one line}
\newcommand{\step}[4]{\vspace{6pt}\begin{center}\includegraphics[height=#3,keepaspectratio]{figs/#2}\end{center}\vspace{-2pt}{\large\textcolor{red}{\textbf{#1}}\ \ #4}\par\vspace{8pt}}
\newcommand{\steptext}[2]{\vspace{4pt}{\large\textcolor{red}{\textbf{#1}}\ \ #2}\par\vspace{6pt}}
\newcommand{\head}[2]{{\Large\bfseries #1}\hfill{\large #2}\par\vspace{4pt}\hrule\vspace{4pt}}
\begin{document}
% ---------------- English ----------------
\head{SeeSound: how to record (teachers)}{English}
\steptext{1}{Open \texttt{http://localhost:8502/?page=collect} in Chrome (on the computer where the program runs).}
\step{2}{s1_setup.png}{9.2cm}{Fill in the form with a code, not a name (red 1--5), and tick both consent boxes (6--7).}
\step{3}{s2_start.png}{2.8cm}{Press \textbf{Start}.}
\newpage
\step{4}{s3_record.png}{10.2cm}{Press the \textbf{blue microphone button}, say the word \textbf{once}, and wait.}
\step{5}{s4_saved.png}{10.2cm}{``Saved ✓'' and the next word appear; cough or wrong word? Press \textbf{Undo} (red 1) and say it again.}
\newpage
\step{6}{s5_refused.png}{5.8cm}{``Too quiet / too loud / too noisy''? Press \textbf{Try again}.}
\step{7}{s6_done.png}{3.8cm}{``All done'': press \textbf{Download}, then \textbf{Next participant}.}
\step{8}{s7_rate.png}{6.6cm}{Rating (alone, later, at \texttt{http://localhost:8501}, tab \textbf{Rate}): listen, then choose Correct / Almost / Not yet for the \textbf{first sound only}.}
\vspace{4pt}\hrule\vspace{4pt}
{\large\textbf{Rules:} do not say the word first; do not correct or praise the sound; record again only for a technical problem.}
\newpage
% ---------------- Traditional Chinese (Taiwan) ----------------
\head{SeeSound 錄音操作指南（老師用）}{繁體中文}
\steptext{1}{在執行程式的電腦上，用 Chrome 開啟 \texttt{http://localhost:8502/?page=collect}。}
\step{2}{s1_setup.png}{9.2cm}{填寫表單，用代碼，不填姓名（紅色 1--5），並勾選兩個同意框（6--7）。}
\step{3}{s2_start.png}{2.8cm}{按\textbf{「開始」}。}
\newpage
\step{4}{s3_record.png}{10.2cm}{按\textbf{藍色麥克風鍵}，念\textbf{一次}單字，然後等待。}
\step{5}{s4_saved.png}{10.2cm}{出現「已儲存 ✓」並換下一個字；咳嗽或念錯字？按\textbf{「取消上一筆錄音」}（紅色 1）再念一次。}
\newpage
\step{6}{s5_refused.png}{5.8cm}{顯示「太小聲／太大聲／太吵」？按\textbf{「再試一次」}。}
\step{7}{s6_done.png}{3.8cm}{出現「完成」：按\textbf{「下載」}，再按\textbf{「下一位」}。}
\step{8}{s7_rate.png}{6.6cm}{評分（之後單獨進行，網址 \texttt{http://localhost:8501}，分頁 \textbf{Rate}）：聽錄音，只評\textbf{第一個子音}：正確／差不多／還沒。}
\vspace{4pt}\hrule\vspace{4pt}
{\large\textbf{規則：}不要先念出單字；不要糾正或稱讚發音；只有錄音出問題才重錄。}
\end{document}
